% ============================================================================
%  Analizador Trifásico IoT — Nota de aplicación
%  Documento: AN-ATI-001 · Rev. 0.1 · Octubre de 2026
%
%  Compilar desde esta carpeta:
%      pdflatex application-note.tex
%      pdflatex application-note.tex
%
%  Requiere una distribución LaTeX con TikZ, booktabs, newtx y longtable.
%  Probado con MiKTeX 25.12 (pdflatex).
% ============================================================================
\documentclass[11pt,a4paper]{article}

\usepackage[utf8]{inputenc}
\usepackage[T1]{fontenc}
\usepackage[spanish,es-nodecimaldot,es-noshorthands]{babel}
\addto\captionsspanish{\renewcommand{\tablename}{Tabla}}
\usepackage{newtxtext,newtxmath}
% La fuente de ancho fijo de newtx no está disponible en todas las
% instalaciones de MiKTeX; se usa Latin Modern Typewriter para \texttt.
\renewcommand{\ttdefault}{lmtt}
\usepackage[a4paper,top=2.6cm,bottom=2.5cm,left=2.3cm,right=2.3cm,headheight=14pt]{geometry}
\usepackage{microtype}
\usepackage{graphicx}
\usepackage{booktabs,tabularx,longtable,array,multirow,makecell}
\usepackage{enumitem}
\usepackage{xcolor}
\usepackage{titlesec}
\usepackage{fancyhdr}
\usepackage{caption}
\usepackage{float}
\usepackage{placeins}
\usepackage{needspace}
\usepackage{amsmath}
\usepackage{tikz}
% Signo de porcentaje robusto: evita el macro \% de babel-español,
% que falla dentro de tabularx/longtable ("Incompatible glue units").
\newcommand{\pct}{\text{\char37}}
\usetikzlibrary{positioning,arrows.meta,fit,backgrounds,calc,shapes.geometric}
\usepackage{hyperref}
\usepackage{url}
\usepackage{ragged2e}
\usepackage{setspace}

% ----------------------------------------------------------------------------
%  Identidad visual
% ----------------------------------------------------------------------------
\definecolor{unalblue}{HTML}{1F4E79}
\definecolor{unalgray}{HTML}{595959}
\definecolor{unalrule}{HTML}{2E75B6}

\hypersetup{
  colorlinks=true,
  linkcolor=unalblue,
  citecolor=unalblue,
  urlcolor=unalrule,
  pdftitle={Analizador Trifásico IoT — Nota de aplicación},
  pdfauthor={Equipo del proyecto Analizador Trifásico IoT},
  pdfsubject={Arquitectura de hardware, protección de entrada y cadena de medida},
  pdfkeywords={medición trifásica, RN8302B, ESP32-S3, ISO7741, Wi-Fi HaLow, RS-485, supercapacitor, medición de energía}
}

% ----------------------------------------------------------------------------
%  Estilo de secciones y página
% ----------------------------------------------------------------------------
\titleformat{\section}
  {\sffamily\Large\bfseries\color{unalblue}}
  {\thesection}{0.8em}{}
\titleformat{\subsection}
  {\sffamily\large\bfseries\color{unalblue!90!black}}
  {\thesubsection}{0.6em}{}
\titleformat{\subsubsection}
  {\sffamily\normalsize\bfseries}
  {\thesubsubsection}{0.5em}{}
% Las secciones sin número (Resumen, Referencias) conservan la identidad visual
\titleformat{name=\section,numberless}
  {\sffamily\Large\bfseries\color{unalblue}}
  {}{0em}{}
\titlespacing*{\section}{0pt}{2.4ex plus .4ex minus .2ex}{1.0ex}
\titlespacing*{\subsection}{0pt}{1.8ex plus .3ex minus .2ex}{0.7ex}

\pagestyle{fancy}
\fancyhf{}
\fancyhead[L]{\sffamily\footnotesize\color{unalgray}AN-ATI-001 · Analizador Trifásico IoT}
\fancyhead[R]{\sffamily\footnotesize\color{unalgray}Rev. 0.1 · Octubre de 2026}
\fancyfoot[C]{\sffamily\footnotesize\thepage}
\renewcommand{\headrulewidth}{0.4pt}
\renewcommand{\footrulewidth}{0pt}

\captionsetup{
  font=small,
  labelfont={bf,color=unalblue},
  justification=centering,
  skip=6pt
}
\setlength{\parindent}{0pt}
\setlength{\parskip}{0.55em}
\renewcommand{\arraystretch}{1.12}
\setlist{nosep,leftmargin=1.4em}

% ----------------------------------------------------------------------------
%  Portada
% ----------------------------------------------------------------------------
\begin{document}

\begin{titlepage}
\thispagestyle{empty}
\begin{flushleft}
{\sffamily\small\color{unalgray}UNIVERSIDAD NACIONAL DE COLOMBIA}\\[2pt]
{\sffamily\small\color{unalgray}FACULTAD DE INGENIERÍA}\\[2pt]
{\sffamily\small\color{unalgray}PROYECTO ANALIZADOR TRIFÁSICO IoT}
\end{flushleft}

\vspace{1.1cm}
{\color{unalblue}\rule{\textwidth}{2.5pt}}
\vspace{0.9cm}

{\sffamily\Huge\bfseries\color{unalblue} Analizador Trifásico IoT\par}
\vspace{0.55cm}
{\sffamily\LARGE Arquitectura de hardware,\\ protección de entrada y cadena de medida\par}
\vspace{0.9cm}
{\sffamily\large\bfseries Nota de aplicación\par}
{\sffamily\large AN-ATI-001 · Rev. 0.1\par}
\vspace{0.35cm}
{\color{unalblue}\rule{\textwidth}{0.8pt}}

\vfill

\begin{tabularx}{\textwidth}{@{}>{\bfseries}p{4.3cm}X@{}}
Documento & Nota de aplicación técnica \\
Versión & 0.1 (borrador) \\
Fecha & Octubre de 2026 \\
Autor & Equipo del proyecto Analizador Trifásico IoT \\
Institución & Universidad Nacional de Colombia \\
Referencia & Arquitectura V2 del proyecto (07/10/2026) \\
Alcance & Hardware; no incluye firmware ni ensayos \\
\end{tabularx}

\vspace{0.9cm}
{\footnotesize\color{unalgray}
Documento de trabajo académico. Describe la arquitectura de hardware y las
decisiones de diseño del instrumento. Los puntos que aún no están definidos se
reúnen en la Sección~\ref{sec:desarrollo}.}
\end{titlepage}

% ----------------------------------------------------------------------------
%  Resumen
% ----------------------------------------------------------------------------
\pagenumbering{roman}
\setcounter{page}{1}

\section*{Resumen}
\addcontentsline{toc}{section}{Resumen}

Este documento presenta la arquitectura de hardware del \textbf{Analizador
Trifásico IoT}, un instrumento académico para la medición de variables
eléctricas en redes trifásicas de baja tensión. La nota se organiza alrededor
de la \textbf{arquitectura V2 del proyecto (07/10/2026)}, adoptada como diseño
vigente.

Se describen la entrada trifásica y su protección, el acondicionamiento
analógico, el AFE metrológico RN8302B, el aislamiento digital con ISO7741, el
módulo ESP32-S3-WROOM-1-N4, la arquitectura de alimentación con respaldo de
supercapacitores y las interfaces RS-485, Wi-Fi HaLow, microSD y USB-C. Se
identifican además los puntos técnicos que deben resolverse antes de fabricar
la PCB; se reúnen en una sola sección para no interrumpir la descripción del
diseño.

\vspace{0.4em}
\noindent\textbf{Palabras clave:} medición trifásica; RN8302B; ESP32-S3;
ISO7741; Wi-Fi HaLow; RS-485; supercapacitor; medición de energía.

\vspace{0.6em}
\noindent\textbf{Alcance.} Esta nota documenta hardware. No especifica
firmware, calibración metrológica, ensayos de conformidad ni el procedimiento
de instalación; esos temas se mencionan solo cuando condicionan una decisión de
hardware.

% ----------------------------------------------------------------------------
%  Índices
% ----------------------------------------------------------------------------
\newpage
\tableofcontents
\newpage
\listoffigures
\vspace{0.5em}
\listoftables

\newpage
\pagenumbering{arabic}
\setcounter{page}{1}

% ----------------------------------------------------------------------------
%  Contenido
% ----------------------------------------------------------------------------
\section{Introducción}

\subsection{Propósito y alcance}

El propósito de esta nota es documentar, con nivel de detalle de diseño, cómo
está construido el hardware del Analizador Trifásico IoT y por qué. El
documento se dirige a las personas que revisarán, replicarán o continuarán el
diseño; por eso cada afirmación sobre componentes se apoya en los datasheets
del fabricante y en la arquitectura de referencia del proyecto.

La nota se limita a los bloques que definen la arquitectura vigente. No es un
manual de usuario, no describe firmware y no presenta resultados de ensayos.
Los criterios de selección, las interfaces y las decisiones de protección se
explican en el cuerpo del documento; los puntos que aún no están definidos se
reúnen en la Sección~\ref{sec:desarrollo} para no interrumpir la descripción
del diseño.

\subsection{Organización del documento}

La nota sigue el recorrido de la señal y de la energía por el instrumento:

\begin{itemize}
  \item La Sección~\ref{sec:sistema} presenta la arquitectura general, el
        principio de operación y los requisitos de producto que condicionan el
        hardware.
  \item La Sección~\ref{sec:entrada} describe la entrada trifásica, los
        fusibles, los varistores y la distribución de las fases.
  \item La Sección~\ref{sec:medicion} detalla los divisores de tensión, las
        entradas de corriente y el AFE metrológico RN8302B.
  \item La Sección~\ref{sec:aislamiento} explica la barrera galvánica y el
        enlace SPI aislado.
  \item La Sección~\ref{sec:procesamiento} cubre el ESP32-S3 y las interfaces
        RS-485, Wi-Fi HaLow, microSD y USB-C.
  \item La Sección~\ref{sec:alimentacion} describe los rieles, el respaldo de
        supercapacitores y el presupuesto de corriente.
  \item La Sección~\ref{sec:desarrollo} reúne los aspectos que deben
        completarse antes de fabricar la PCB.
\end{itemize}

\subsection{Convenciones y nomenclatura}

\textbf{Alcance de la revisión.} Revisión 0.1, borrador técnico. Describe la
arquitectura de hardware y las decisiones tomadas; el prototipo no se ha
fabricado. Los valores que dependen de una selección posterior se tratan en la
Sección~\ref{sec:desarrollo}.

\textbf{Convenciones.} Las tensiones son eficaces salvo que se indique otra
cosa. Las corrientes son de continua o de pico según el contexto. ``3P4W''
designa una red de tres fases y neutro; ``3P3W'', una red de tres fases sin
neutro.

\begin{table}[H]
\centering
\caption{Abreviaturas usadas en el documento.}
\label{tab:abbr}
\small
\begin{tabularx}{\textwidth}{@{}>{\bfseries}p{2.0cm} X@{}}
\toprule
\textbf{Término} & \textbf{Significado} \\
\midrule
AFE & \emph{Analog front end}; cadena de acondicionamiento y conversión analógica. \\
CT & Transformador de corriente (\emph{current transformer}). \\
MOV & Varistor de óxido metálico (\emph{metal-oxide varistor}). \\
SPI & \emph{Serial Peripheral Interface}. \\
GPIO & \emph{General-purpose input/output}. \\
UART & \emph{Universal asynchronous receiver-transmitter}. \\
SDIO & \emph{Secure Digital Input/Output}. \\
JTAG & \emph{Joint Test Action Group}; interfaz de programación y depuración. \\
BOM & \emph{Bill of materials}; lista de materiales. \\
Last Gasp & Operación de respaldo para cerrar el registro de forma ordenada tras un corte de red. \\
\bottomrule
\end{tabularx}
\end{table}
\Needspace{5\baselineskip}
\section{Descripción del sistema}
\label{sec:sistema}

\subsection{Diagrama de bloques}

La Figura~\ref{fig:bloques} presenta el diagrama de bloques del instrumento,
con los bloques, las interfaces y los rangos de operación del diseño. Las
líneas azules representan señales; las verdes, distribución de energía.

\begin{figure}[H]
\centering
\begin{tikzpicture}[
  x=1cm,y=1cm,
  block/.style={draw, rounded corners=2pt, align=center,
    font=\sffamily\scriptsize, inner sep=3pt, minimum height=1.05cm},
  input/.style={block, fill=red!8, draw=red!55!black,
    text width=3.25cm, minimum width=3.45cm},
  prot/.style={block, fill=orange!10, draw=orange!55!black,
    text width=3.25cm, minimum width=3.45cm},
  meas/.style={block, fill=yellow!18, draw=yellow!45!black,
    text width=3.25cm, minimum width=3.45cm},
  iso/.style={block, fill=purple!8, draw=purple!45!black,
    text width=3.25cm, minimum width=3.45cm},
  dig/.style={block, fill=blue!8, draw=blue!45!black,
    text width=3.25cm, minimum width=3.45cm},
  pwr/.style={block, fill=green!8, draw=green!45!black,
    text width=2.45cm, minimum width=2.65cm},
  per/.style={block, fill=cyan!8, draw=cyan!45!black,
    text width=2.35cm, minimum width=2.55cm},
  sig/.style={-{Latex[length=1.7mm]}, semithick, draw=blue!55!black,
    rounded corners=8pt},
  sigbus/.style={semithick, draw=blue!55!black, rounded corners=8pt},
  pwra/.style={-{Latex[length=1.7mm]}, semithick, draw=green!45!black,
    rounded corners=8pt},
  pwrbus/.style={semithick, draw=green!45!black, rounded corners=8pt},
  lab/.style={font=\sffamily\tiny, fill=white, inner sep=1.2pt},
  stage/.style={font=\sffamily\tiny\bfseries, text=black!60}
]

% Encabezados funcionales
\node[stage] at (1.5,1.35) {ALIMENTACIÓN};
\node[stage] at (7.4,1.35) {MEDIDA, AISLAMIENTO Y CONTROL};

% Columna izquierda: generación, respaldo y rieles
\node[pwr] (ps1) at (0,0)
  {PS1 · HLK-10M05\\5 V / 2 A};
\node[pwr] (ps2) at (3.0,0)
  {PS2 · HLK-5M03B\\AVCC\_3v3 / 1,5 A};
\node[pwr] (backup) at (0,-2.0)
  {Respaldo Last Gasp\\SY6280AAC\\2$\times$30 F en serie};
\node[pwr] (buck) at (0,-4.0)
  {U3 · AP63203WU-7\\DVCC\_3v3 / 2 A};

% Columna central: flujo principal inspirado en el diagrama V2
\node[input] (mains) at (7.4,0)
  {Red trifásica\\L1 / L2 / L3 / N};
\node[prot] (term) at (7.4,-1.55)
  {Bornera U4 · 4 vías\\C50126V-4P0G56};
\node[prot] (prot) at (7.4,-3.1)
  {Protección de entrada\\F1--F3 (2 A) + RV1--RV3};
\node[meas] (cond) at (7.4,-4.65)
  {Acondicionamiento analógico\\divisores + filtros RC};
\node[meas] (afe) at (7.4,-6.2)
  {AFE RN8302B\\3 V + 4 I};
\node[iso] (iso) at (7.4,-7.75)
  {ISO7741 · 4 canales\\aislamiento SPI};
\node[dig] (mcu) at (7.4,-9.3)
  {\mbox{ESP32-S3-WROOM-1-N4}\\procesamiento y control};

% Programación a la derecha y periféricos bajo el procesador
\node[per] (usb) at (12.0,-9.3)
  {USB-C · J401\\programación / monitoreo};
\node[per] (rs485) at (4.0,-11.8)
  {RS-485\\VISOM RSM3485PCT};
\node[per] (halow) at (7.4,-11.8)
  {Wi-Fi HaLow\\Quectel FGH100M};
\node[per] (sd) at (10.8,-11.8)
  {microSD\\tf-card-9p};

% Flujo vertical de entrada, medida y control
\draw[sig] (mains) -- (term);
\draw[sig] (term) -- (prot);
\draw[sig] (prot) -- node[lab,right=3pt] {Phase\_1/2/3/N} (cond);
\draw[sig] (cond) -- (afe);
\draw[sig] (afe) -- node[lab,right=3pt] {SPI} (iso);
\draw[sig] (iso) -- node[lab,right=3pt] {SPI aislado} (mcu);
\draw[sig] (mcu) -- node[lab,above=2pt] {USB / JTAG} (usb);

% Derivación protegida hacia las dos fuentes, sin cruzar sus salidas
\draw[sigbus] (prot.west) -- (4.85,-3.1) -- (4.85,0.82) -- (0,0.82);
\node[lab] at (2.6,0.98) {Phase\_1 / N};
\draw[sig] (0,0.82) -- (ps1.north);
\draw[sig] (3.0,0.82) -- (ps2.north);

% Rama de respaldo y fuente auxiliar de medida
\draw[pwra] (ps1) -- (backup);
\draw[pwra] (backup) -- (buck);
\draw[pwra] (ps2.south) -- (3.0,-6.2)
  node[lab,pos=0.72,right=3pt] {AVCC\_3v3} -- (afe.west);
\node[lab] at (4.35,-6.02) {10 mA};

% Distribución DVCC: bus lateral independiente de las señales
\draw[pwrbus] (buck.west) -- (-1.65,-4.0) -- (-1.65,-13.25)
  -- (10.8,-13.25);
\draw[pwra] (-1.65,-9.3) -- node[lab,above=2pt]
  {DVCC\_3v3 · 500 mA} (mcu.west);
\draw[pwra] (4.0,-13.25) -- node[lab,pos=0.55,right=3pt]
  {90 mA} (rs485.south);
\draw[pwra] (7.4,-13.25) -- node[lab,pos=0.55,right=3pt]
  {350 mA} (halow.south);
\draw[pwra] (10.8,-13.25) -- node[lab,pos=0.55,right=3pt]
  {200 mA} (sd.south);

% Interfaces de los periféricos: ramas paralelas y equidistantes
\draw[sigbus] (mcu.south) -- (7.4,-10.72) -- (4.0,-10.72);
\draw[sigbus] (7.4,-10.72) -- (10.8,-10.72);
\draw[sig] (4.0,-10.72) -- node[lab,pos=0.55,left=3pt]
  {UART} (rs485.north);
\draw[sig] (7.4,-10.72) -- node[lab,pos=0.55,left=3pt]
  {SDIO / SPI} (halow.north);
\draw[sig] (10.8,-10.72) -- node[lab,pos=0.55,left=3pt]
  {SPI} (sd.north);

\end{tikzpicture}
\caption{Diagrama de bloques del Analizador Trifásico IoT. Las líneas azules
representan señales y las verdes, distribución de energía.}
\label{fig:bloques}
\end{figure}

\subsection{Principio de operación}

La cadena de medida comienza en la bornera de entrada U4, donde se conectan
L1, L2, L3 y N. Cada fase pasa por un fusible y un varistor antes de llegar a
las etapas de potencia y de medida. En la etapa de medida, las tensiones se
atenúan con divisores resistivos y las corrientes entran por cuatro conectores
para transformadores de corriente. El RN8302B convierte y procesa las señales;
el ISO7741 aísla el bus SPI que las lleva al ESP32-S3. El microcontrolador
agrega la energía, sella la hora, escribe en la microSD y atiende las
comunicaciones RS-485 y Wi-Fi HaLow.

La alimentación se deriva de \texttt{Phase\_1} y \texttt{N}. La fuente PS1
(HLK-10M05) entrega 5 V y carga dos supercapacitores de 30 F en serie; el buck
U3 (AP63203WU-7) genera el riel digital de 3.3 V. La fuente PS2 (HLK-5M03B)
genera el riel auxiliar de 3.3 V para el dominio de medida. El respaldo de
supercapacitores permite el cierre ordenado del registro tras un corte de red
(\emph{Last Gasp}).

\subsection{Requisitos de referencia}

La Tabla~\ref{tab:requisitos} resume los requisitos de producto que condicionan
el hardware \cite{rdem01}. La última columna indica cómo los cubre el diseño
actual o qué falta para cubrirlos. No se presentan como resultados de ensayo.

\begin{table}[H]
\centering
\caption{Requisitos de producto relevantes para el hardware.}
\label{tab:requisitos}
\footnotesize
\begin{tabularx}{\textwidth}{@{}>{\bfseries}p{1.0cm} >{\raggedright\arraybackslash}p{6.3cm} X@{}}
\toprule
\textbf{ID} & \textbf{Requisito} & \textbf{Cobertura en el diseño} \\
\midrule
F-01 & Medición de V, I, P, Q, S, PF y frecuencia, por fase y total, en 3P4W y 3P3W. & AFE RN8302B y firmware; el hardware cubre 3P4W. El conexionado 3P3W debe verificarse. \\
F-03 & Error $\leq 2\,\pct$ en energía activa (clase 2), de 5\,\pct a 120\,\pct de $I_n$. & Objetivo de exactitud; requiere calibración. \\
F-04 & 90--264 V por canal de tensión, 50/60 Hz, 1,2 $U_n$ permanente. & Parcial; verificar divisores, bornera y fusibles para 440 V entre fases. \\
F-05 & 0,05--100 A por fase con TC de núcleo partido de 333 mV. & Parcial; modelo de TC por definir. \\
F-06 & Alimentación auxiliar de 85--264 VCA, independiente de la red medida. & Parcial; PS1 y PS2 se alimentan de \texttt{Phase\_1}/\texttt{N}. \\
F-07 & Modbus RTU esclavo por RS-485, 9,6--115,2 kbps; Wi-Fi HaLow con MQTT/TLS. & RS-485 implementado; integración de HaLow por completar. \\
F-08 & Registro en microSD de al menos 90 días a 15 minutos, CSV con verificación de integridad. & microSD implementada; registro a cargo del firmware. \\
F-09 & Al menos 5 s de operación tras un corte, con cierre ordenado. & Supercapacitores dimensionados; falta ensayo. \\
NF-02 & Aislamiento de 2,5 kVCA entre medida y comunicación/USB. & ISO7741DWR de 5 kVrms; revisar separación de tierras (Sección~\ref{sec:aislamiento}). \\
NF-03 & Inmunidad IEC 61000-4-2 nivel 3, 4-4 nivel 3, 4-5 nivel 4 y 4-6 nivel 3; emisión EN 55011 clase B. & Requiere ensayo y layout definitivo. \\
NF-10 & Categoría de medida III, 300 V; IEC 61010-1. & Requiere revisión normativa. \\
\bottomrule
\end{tabularx}
\end{table}
\section{Entrada trifásica y protección}
\label{sec:entrada}

\subsection{Topología y conexionado}

La entrada se realiza por \textbf{U4}, una bornera de cuatro vías de 5 mm
(JXTCONN C50126V-4P0G56) \cite{jxtconn}. Las cuatro vías corresponden a L1, L2,
L3 y N. Cada fase pasa por un fusible de 2 A (F1--F3, JDT Fuse
JFC2410-0200TS, formato 2410, \emph{time-lag}) \cite{jfcfuse}. El neutro llega
directo a la red \texttt{N}. Después de los fusibles, cada fase tiene un varistor
MOV-14D471K (RV1--RV3) conectado a la referencia de protección \texttt{GND}.
Las fases protegidas salen como \texttt{Phase\_1}, \texttt{Phase\_2},
\texttt{Phase\_3} y \texttt{N} hacia las etapas de potencia y de medida. La
Figura~\ref{fig:entrada} dibuja la etapa como esquema eléctrico: los fusibles y
los varistores aparecen como componentes, no como bloques, y las uniones se
trazan como cables; el arco sobre un cruce indica que los conductores no se
conectan.

\begin{figure}[H]
\centering
\begin{tikzpicture}[
  x=1cm,y=1cm,
  block/.style={draw,rounded corners=2pt,align=center,font=\sffamily\scriptsize,inner sep=3pt,minimum height=0.9cm},
  port/.style={block,fill=blue!8,draw=blue!55!black,minimum width=2.0cm},
  wire/.style={semithick,draw=blue!55!black},
  lab/.style={font=\sffamily\scriptsize},
  desig/.style={font=\sffamily\scriptsize,inner sep=1pt}
]
% U4: bornera de entrada
\node[block,fill=gray!8,draw=gray!55!black,minimum width=2.2cm,minimum height=4.3cm]
  (u4) at (0,0) {U4\\Bornera\\4 vías};
\node[lab,anchor=east] at (-1.15,1.6) {L1};
\node[lab,anchor=east] at (-1.15,0.4) {L2};
\node[lab,anchor=east] at (-1.15,-0.8) {L3};
\node[lab,anchor=east] at (-1.15,-2.0) {N};

% Puertos jerárquicos de salida
\node[port] (p1) at (7.6,1.6) {Phase\_1};
\node[port] (p2) at (7.6,0.4) {Phase\_2};
\node[port] (p3) at (7.6,-0.8) {Phase\_3};
\node[port] (pn) at (7.6,-2.0) {N};

% Cables de fase y de neutro
\draw[wire] (1.1,1.6) -- (6.6,1.6);
\draw[wire] (1.1,0.4) -- (6.6,0.4);
\draw[wire] (1.1,-0.8) -- (6.6,-0.8);
\draw[wire] (1.1,-2.0) -- (6.6,-2.0);

% Fusibles F1--F3 (símbolo IEC: rectángulo atravesado por el conductor)
\foreach \y/\f in {1.6/F1, 0.4/F2, -0.8/F3}{
  \draw[wire] (2.58,\y+0.11) rectangle (3.12,\y-0.11);
  \node[desig,anchor=south] at (2.85,\y+0.15) {\f};
  \node[desig,anchor=north] at (2.85,\y-0.15) {2\,A};
}

% Nodos de derivación hacia los varistores
\fill[blue!55!black] (4.2,1.6) circle (0.035);
\fill[blue!55!black] (5.1,0.4) circle (0.035);
\fill[blue!55!black] (6.0,-0.8) circle (0.035);

% Varistores RV1--RV3: cada fase deriva a GND.
% El arco en un cruce indica conductores que se cruzan sin conexión.
\draw[wire] (4.2,1.6) -- (4.2,0.48) arc (90:-90:0.08)
  -- (4.2,-0.72) arc (90:-90:0.08) -- (4.2,-1.92) arc (90:-90:0.08)
  -- (4.2,-2.85);
\draw[wire] (5.1,0.4) -- (5.1,-0.72) arc (90:-90:0.08)
  -- (5.1,-1.92) arc (90:-90:0.08) -- (5.1,-2.85);
\draw[wire] (6.0,-0.8) -- (6.0,-1.92) arc (90:-90:0.08) -- (6.0,-2.85);

\foreach \x/\r in {4.2/RV1, 5.1/RV2, 6.0/RV3}{
  \draw[wire] (\x-0.11,-3.35) rectangle (\x+0.11,-2.85);
  \draw[wire] (\x-0.15,-3.39) -- (\x+0.15,-2.81);
  \draw[wire] (\x,-3.35) -- (\x,-3.60);
  \node[desig,anchor=south west] at (\x+0.14,-2.79) {\r};
}

% Riel de referencia de protección
\draw[thick] (4.2,-3.60) -- (6.6,-3.60);
\draw[thick] (6.6,-3.60) -- (6.6,-3.73);
\draw[thick] (6.38,-3.73) -- (6.82,-3.73);
\draw[thick] (6.46,-3.83) -- (6.74,-3.83);
\draw[thick] (6.53,-3.93) -- (6.67,-3.93);
\node[lab,anchor=north] at (6.6,-4.01) {GND (referencia de protección)};

\end{tikzpicture}
\caption{Cadena de entrada y protección, dibujada como esquema eléctrico. Cada
fase pasa por su fusible F1--F3 (2\,A) y su varistor RV1--RV3 deriva a la
referencia \texttt{GND}; el neutro llega directo a la etapa de medida. El arco
sobre un cruce indica que los conductores no se conectan.}
\label{fig:entrada}
\end{figure}

\subsection{Componentes de entrada}

\begin{table}[H]
\centering
\caption{Componentes de la etapa de entrada y protección.}
\label{tab:entrada}
\footnotesize
\begin{tabularx}{\textwidth}{@{}>{\bfseries}p{1.5cm} >{\raggedright\arraybackslash}p{3.2cm} >{\raggedright\arraybackslash}X >{\raggedright\arraybackslash}p{4.3cm}@{}}
\toprule
\textbf{Desig.} & \textbf{Componente} & \textbf{Valor / referencia} & \textbf{Función y destino} \\
\midrule
U4 & Bornera de 4 vías, 5 mm & C50126V-4P0G56, 18 A, 250 V & Entrada L1/L2/L3/N; alimenta F1--F3 y \texttt{N}. \\
F1--F3 & Fusible 2410, \emph{time-lag} & JFC2410-0200TS, 2 A & Protección de cada fase; salida \texttt{Phase\_1/2/3}. \\
RV1--RV3 & Varistor de óxido metálico, 14 mm & MOV-14D471K, 470 V & Supresión de transitorios fase--\texttt{GND}. \\
U101--U104 & Conector de 3 pines & HX PJ-320A-3P DIP & Entrada de los cuatro canales de corriente. \\
\bottomrule
\end{tabularx}
\end{table}

\subsection{Distribución de las fases}

Las tres fases protegidas y el neutro se distribuyen con los nombres
\texttt{Phase\_1}, \texttt{Phase\_2}, \texttt{Phase\_3} y \texttt{N}. La etapa
de medida recibe las cuatro señales; la etapa de potencia toma únicamente
\texttt{Phase\_1} y \texttt{N}. En consecuencia:

\begin{itemize}
  \item La medición de tensión recibe las tres fases y el neutro.
  \item La alimentación se toma de una sola fase (\texttt{Phase\_1}) y del
        neutro, a través de las ferritas FB1 y FB2.
  \item El neutro no lleva fusible; esta decisión debe revisarse junto con la
        coordinación de protecciones.
\end{itemize}

\subsection{Observaciones de protección}

\begin{itemize}
  \item \textbf{Tensión de la bornera.} La ficha del proveedor de U4 indica
        250 V. El requisito de producto contempla 440 V entre fases en 3P3W;
        la tensión nominal de la bornera queda por debajo. Debe verificarse y,
        si aplica, seleccionar una bornera de mayor tensión.
  \item \textbf{Tensión de los fusibles.} Los fusibles 2410 se especifican
        típicamente para 250 VCA. En 3P4W cada fusible ve la tensión
        fase--neutro; en 3P3W ve la tensión entre fases. La coordinación debe
        hacerse para el caso más exigente.
  \item \textbf{Coordinación fusible--varistor--fuente.} El fabricante del
        HLK-5M03 recomienda un fusible de 1 A/250 VCA para la fuente; el diseño
        usa 2 A. La selección debe justificarse con la corriente de entrada, la
        corriente de irrupción y la capacidad de la fuente.
  \item \textbf{Referencia de protección.} RV1--RV3 van de fase a \texttt{GND}.
        La unión de \texttt{GND} con la tierra de protección (PE) y el
        comportamiento ante una falla de fase a tierra deben definirse antes de
        fabricar.
\end{itemize}
\section{Acondicionamiento analógico y AFE}
\label{sec:medicion}

\subsection{Canales de tensión}

El frontal de medida tiene tres canales diferenciales de tensión. La red de
atenuación está formada por resistencias de precisión de
$165\,\mathrm{k\Omega}$ (R1--R6, R9--R14 y R17--R22) y resistencias de
$1\,\mathrm{k\Omega}$ (R7, R8, R15, R16, R23 y R24), con capacitores de
$33\,\mathrm{nF}$ (C2, C10, C12, C13, C14 y C15) para el filtrado. Cada fase
usa seis resistencias de $165\,\mathrm{k\Omega}$ y dos de $1\,\mathrm{k\Omega}$.

El rango de entrada de los ADC de tensión se fija en \textbf{800 mVpp}. El
divisor se dimensiona para que la tensión máxima de fase produzca ese rango en
la entrada del AFE.

\subsection{Canales de corriente}

Hay cuatro canales diferenciales de corriente: tres fases y neutro. Cada canal
entra por un conector de tres pines HX PJ-320A-3P DIP (U101--U104) y pasa por
una red RC formada por resistencias de $1\,\mathrm{k\Omega}$ (R101--R108) y
capacitores de $33\,\mathrm{nF}$ (C101--C108). El rango de entrada es de
\textbf{50 mVpp a 400 mVpp} y el punto nominal de operación es de
$\approx 150\,\mathrm{mVpp}$. El rango es adaptativo y el punto de operación
depende del transformador de corriente seleccionado.

El diseño prevé un TC de núcleo partido con salida de 333 mV y un rango de
0,05 A a 100 A.

\subsection{AFE metrológico RN8302B}

El AFE es el \textbf{RENERGY RN8302B} (U2). Integra tres canales de tensión y
cuatro de corriente, y calcula internamente tensión, corriente, potencia activa,
reactiva y aparente, factor de potencia y distorsión armónica. El
microcontrolador no realiza la metrología; lee los registros del RN8302B por
SPI.

El reloj de metrología es un cristal dedicado \textbf{X101}
(\texttt{049ss2-8.192f12dzlal}) de \textbf{8,192 MHz}, con dos capacitores de
carga de $15\,\mathrm{pF}$ (C1 y C11). El desacople del AFE comprende cuatro
capacitores de $100\,\mathrm{nF}$, dos de $10\,\mathrm{\mu F}$ y uno de
$1\,\mathrm{\mu F}$, que corresponden a C4, C6, C8 y C9; C5 y C7; y C3,
respectivamente.

\begin{table}[H]
\centering
\caption{Componentes principales de la cadena de medida.}
\label{tab:medicion}
\footnotesize
\begin{tabularx}{\textwidth}{@{}>{\raggedright\arraybackslash}p{3.6cm} >{\raggedright\arraybackslash}p{5.3cm} X@{}}
\toprule
\textbf{Sub-bloque} & \textbf{Componentes} & \textbf{Función} \\
\midrule
Divisores de tensión &
R1--R6, R9--R14, R17--R22 ($165\,\mathrm{k\Omega}$);
R7, R8, R15, R16, R23, R24 ($1\,\mathrm{k\Omega}$);
C2, C10, C12--C15 ($33\,\mathrm{nF}$) &
Atenuación y filtrado de las tres fases. \\
Entradas de corriente &
U101--U104 (HX PJ-320A-3P DIP);
R101--R108 ($1\,\mathrm{k\Omega}$);
C101--C108 ($33\,\mathrm{nF}$) &
Conector y filtro RC de cada canal de corriente. \\
AFE &
U2 (RN8302B);
X101 (8,192 MHz);
C1, C11 ($15\,\mathrm{pF}$);
C3--C9 (desacople) &
Conversión, cálculo metrológico y reloj. \\
Aislamiento &
U1 (ISO7741DWR);
R25, R26 ($10\,\mathrm{k\Omega}$) &
Enlace SPI aislado; pull-ups de EN1/EN2. \\
\bottomrule
\end{tabularx}
\end{table}

\subsection{Filtrado y disposición}

Las redes RC de entrada atenúan el ruido de alta frecuencia antes del AFE. Los
capacitores de desacople se ubican junto a los pines de alimentación del
circuito integrado y las pistas diferenciales se trazan simétricas y cortas
para conservar la simetría entre canales.
\section{Aislamiento galvánico}
\label{sec:aislamiento}

\subsection{ISO7741 y señales aisladas}

El aislamiento entre el dominio de medida y el dominio digital se realiza con
\textbf{U1}, un ISO7741DWR de cuatro canales. El ISO7741 tiene tres canales en
un sentido y uno en sentido inverso, lo que coincide con un bus SPI: SCLK, SCS
y SDI van del microcontrolador al AFE, y SDIO regresa del AFE al
microcontrolador.

\begin{figure}[H]
\centering
\begin{tikzpicture}[
  x=1cm,y=1cm,
  dombox/.style={draw, rounded corners=2pt, minimum width=3.55cm,
    minimum height=5.2cm},
  measdom/.style={dombox, fill=yellow!10, draw=yellow!45!black},
  digdom/.style={dombox, fill=blue!6, draw=blue!45!black},
  iso/.style={draw, rounded corners=2pt, fill=purple!8,
    draw=purple!45!black, minimum width=4.9cm, minimum height=5.2cm},
  sig/.style={-{Latex[length=1.8mm]}, semithick, draw=blue!55!black,
    rounded corners=8pt},
  chan/.style={font=\sffamily\scriptsize\bfseries, text=black!80},
  pin/.style={font=\sffamily\tiny, text=black!65},
  domtag/.style={font=\sffamily\tiny\bfseries, text=black!55}
]

% Dominios de referencia: cada uno ocupa un lado de la barrera
\draw[measdom] (-7.1,-2.6) rectangle (-3.35,2.6);
\draw[digdom]   (3.35,-2.6) rectangle (7.1,2.6);

\node[domtag] at (-5.225,2.25) {DOMINIO DE MEDIDA};
\node[font=\sffamily\large\bfseries] at (-5.225,1.15) {RN8302B};
\node[font=\sffamily\scriptsize, text=black!65] at (-5.225,0.62) {AFE};
\node[font=\sffamily\scriptsize] at (-5.225,-1.6) {AVCC\_3v3};
\node[pin] at (-5.225,-2.05) {referencia GND1};

\node[domtag] at (5.225,2.25) {DOMINIO DIGITAL};
\node[font=\sffamily\large\bfseries] at (5.225,1.15) {ESP32-S3};
\node[font=\sffamily\scriptsize, text=black!65] at (5.225,0.62) {procesamiento};
\node[font=\sffamily\scriptsize] at (5.225,-1.6) {DVCC\_3v3};
\node[pin] at (5.225,-2.05) {referencia GND2};

% Aislador: cuatro canales alineados y barrera en el centro del componente
\draw[iso] (-2.45,-2.6) rectangle (2.45,2.6);
\node[font=\sffamily\small\bfseries] at (0,2.25) {ISO7741DWR};
\node[pin] at (0,1.85) {4 canales · aislamiento SPI};

\foreach \yy in {0.7,-0.3,-1.3} {
  \draw[black!10] (-2.3,\yy) -- (-0.15,\yy);
  \draw[black!10] (0.15,\yy) -- (2.3,\yy);
}

\node[chan,anchor=east] at (-0.3,1.2)  {SCLK};
\node[pin,anchor=west]  at (0.3,1.2)   {INC · pin 5};
\node[chan,anchor=east] at (-0.3,0.2)  {SCS};
\node[pin,anchor=west]  at (0.3,0.2)   {INB · pin 4};
\node[chan,anchor=east] at (-0.3,-0.8) {SDI};
\node[pin,anchor=west]  at (0.3,-0.8)  {IND · pin 6};
\node[chan,anchor=east] at (-0.3,-1.8) {SDIO};
\node[pin,anchor=west]  at (0.3,-1.8)  {INA · pin 3};

\node[pin,anchor=east] at (-0.2,-2.3) {VCC1 1 · GND1 2/8};
\node[pin,anchor=west] at (0.2,-2.3)  {VCC2 16 · GND2 9/15};

% Tres canales van del MCU al AFE; SDIO regresa del AFE al MCU
\foreach \yy/\lab in {1.2/SCLK,0.2/SCS,-0.8/SDI} {
  \draw[sig] (3.35,\yy) -- (2.45,\yy);
  \draw[sig] (-2.45,\yy) -- (-3.35,\yy);
}
\draw[sig] (-3.35,-1.8) -- (-2.45,-1.8);
\draw[sig] (2.45,-1.8) -- (3.35,-1.8);

% Barrera física: atraviesa el aislador y separa las dos referencias.
% Se interrumpe en el encabezado del ISO7741 para no cruzar su nombre.
\draw[dashed, very thick, red!60!black] (0,-3.25) -- (0,1.65);
\draw[dashed, very thick, red!60!black] (0,2.6) -- (0,3.25);
\node[font=\sffamily\scriptsize\bfseries, text=red!60!black,
  fill=white, inner sep=2pt] at (0,3.4)
  {BARRERA DE AISLAMIENTO GALVÁNICO};

\end{tikzpicture}
\caption{Esquema del aislamiento SPI. El ISO7741 separa tres señales de ida y
una de retorno; la barrera se materializa con las tierras GND1 y GND2
separadas.}
\label{fig:aislamiento}
\end{figure}

\begin{table}[H]
\centering
\caption{Señales que cruzan la barrera de aislamiento.}
\label{tab:aislamiento}
\small
\begin{tabularx}{\textwidth}{@{}>{\bfseries}p{1.6cm} p{2.1cm} >{\raggedright\arraybackslash}p{2.7cm} p{2.4cm} X@{}}
\toprule
\textbf{Señal} & \textbf{Pin ISO7741} & \textbf{Nombre en el AFE} & \textbf{Dirección} & \textbf{Función} \\
\midrule
SCLK & 5 (INC) & SCLK\_MCU & MCU $\rightarrow$ AFE & Reloj SPI. \\
SCS  & 4 (INB) & SCS\_MCU  & MCU $\rightarrow$ AFE & Selección de chip. \\
SDI  & 6 (IND) & SDI\_MCU  & MCU $\rightarrow$ AFE & Datos MOSI. \\
SDIO & 3 (INA) & SDIO\_MCU & AFE $\rightarrow$ MCU & Datos MISO. \\
\bottomrule
\end{tabularx}
\end{table}

\noindent\textbf{Capacidad de canales.} El ISO7741 aísla cuatro señales. Si más
adelante se agregan dos GPIO de control del AFE, será necesario un segundo
aislador o un dispositivo de más canales.

\subsection{Dominios de alimentación y referencia}

El lado 1 del ISO7741 se alimenta con \texttt{DVCC\_3v3} (pin 1) y el lado 2 con
\texttt{AVCC\_3v3} (pin 16). Las tierras correspondientes son GND1 (pines 2 y
8) y GND2 (pines 9 y 15). Para que la barrera sea efectiva, GND1 y GND2 deben
ser referencias distintas: la medida usa \texttt{AGND} y el dominio digital
\texttt{GND\_DIG}. Las dos referencias solo pueden unirse en un punto definido
si el diseño lo exige, y esa unión debe reflejarse en el layout.

El ISO7741DWR está especificado para 5000 Vrms de aislamiento en el paquete
DW-16 \cite{iso7741}. El requisito del producto es 2,5 kVCA \cite{rdem01}; el
margen del componente es suficiente, pero la separación real depende de las
referencias, del \emph{creepage} y del \emph{clearance} en la PCB.

\subsection{Integridad de la señal y layout}

El enlace SPI se traza con pistas cortas y simétricas, y el retorno de
corriente se mantiene próximo a cada señal. La separación entre la barrera y
los planos de cobre se verifica en el layout, junto con el tratamiento de los
pines EN1/EN2.
\section{Procesamiento y periféricos}
\label{sec:procesamiento}

\subsection{ESP32-S3-WROOM-1-N4}

El procesamiento se realiza con el módulo \textbf{ESP32-S3-WROOM-1-N4} (U5),
que integra el SoC ESP32-S3 y 4 MB de flash. El módulo concentra el bus SPI
aislado del AFE, la UART de RS-485, la interfaz del módulo Wi-Fi HaLow, el bus
SPI de la microSD y el puerto USB-C de programación y monitorización.

Al módulo lo acompañan U6 (RS-485), U7 (microSD) y J401 (USB-C), además de los
capacitores de desacople, las resistencias de pull-up y las protecciones de ESD
que requiere cada interfaz.

\subsection{RS-485}

La interfaz Modbus RTU se implementa con el módulo aislado
\textbf{VISOM RSM3485PCT} (U6). El módulo integra aislamiento de señal y de
alimentación, y convierte los niveles lógicos del ESP32-S3 a niveles
diferenciales RS-485. La interfaz usa UART (TX/RX) más una GPIO de control y
consume 90 mA sobre el riel de 3.3 V.

\subsection{Wi-Fi HaLow}

El enlace inalámbrico usa el módulo \textbf{Quectel FGH100M} (Morse Micro
MM6108), con interfaces de 6 SDIO, 4 SPI, 2 UART y 2 GPIO, y un consumo
estimado de 350 mA. El módulo se integra siguiendo el diseño de referencia del
fabricante, con su propia alimentación y el conexionado de sus interfaces.

\subsection{Memoria microSD}

La tarjeta microSD se conecta mediante el conector \textbf{tf-card-9p} (U7).
La interfaz es SPI de cuatro hilos y el consumo estimado es de 200 mA. El
conector y la tarjeta permiten el registro local de 90 días a intervalos de
15 minutos, con verificación de integridad a cargo del firmware.

\subsection{USB-C y programación}

El conector \textbf{J401} (\texttt{CON\_USB\_TYPEC\_16P}) se usa para
programación y monitorización por USB, en el camino USB-C $\rightarrow$ JTAG del
ESP32-S3. El puerto requiere protecciones ESD y resistencias de terminación si
se expone fuera del gabinete.

\subsection{Mapa de interfaces}

\begin{table}[H]
\centering
\caption{Mapa de interfaces del ESP32-S3-WROOM.}
\label{tab:interfaces}
\small
\begin{tabularx}{\textwidth}{@{}>{\raggedright\arraybackslash}p{4.7cm} >{\raggedright\arraybackslash}p{3.4cm} X@{}}
\toprule
\textbf{Interfaz} & \textbf{Uso} & \textbf{Destino} \\
\midrule
SPI (4 hilos) & Datos y control & RN8302B, a través del ISO7741 \\
UART (TX/RX) + 1 GPIO & Modbus RTU & VISOM RSM3485PCT \\
6 SDIO / 4 SPI / 2 UART / 2 GPIO & Enlace de radio & Quectel FGH100M \\
SPI (4 hilos) & Tarjeta de memoria & microSD \texttt{tf-card-9p} \\
USB (D+/D$-$) & Programación y monitor & Conector USB-C J401 \\
\bottomrule
\end{tabularx}
\end{table}

\noindent\textbf{Asignación de canales.} El bus SPI del AFE ocupa cuatro
señales; si se implementan los dos GPIO de control previstos, requieren un
segundo aislador (Sección~\ref{sec:aislamiento}).
\section{Arquitectura de alimentación}
\label{sec:alimentacion}

\subsection{Rieles y fuentes}

La alimentación se deriva de \texttt{Phase\_1} y \texttt{N}. La fuente
principal \textbf{PS1} (Hi-Link HLK-10M05) entrega 5 V a 2 A y alimenta, a
través del limitador de carga, el respaldo de supercapacitores y el buck
digital. La fuente auxiliar \textbf{PS2} (Hi-Link HLK-5M03B) entrega 3.3 V para
el dominio de medida. El buck \textbf{U3} (AP63203WU-7) genera el riel digital
de 3.3 V a partir de los 5 V.

\begin{table}[H]
\centering
\caption{Rieles de alimentación y fuentes.}
\label{tab:rieles}
\small
\begin{tabularx}{\textwidth}{@{}>{\bfseries}p{2.7cm} >{\raggedright\arraybackslash}p{4.3cm} >{\raggedright\arraybackslash}X@{}}
\toprule
\textbf{Riel} & \textbf{Fuente} & \textbf{Especificación y destino} \\
\midrule
5 V & PS1: HLK-10M05 & 5 V / 2 A, 10 W. Alimenta la entrada del SY6280AAC. \\
5 V limitado & SY6280AAC & Interruptor de distribución con límite de corriente programable. Su salida carga el banco C16/C17 y alimenta la entrada de U3. \\
AVCC\_3v3 & PS2: HLK-5M03B & 3.3 V, 5 W (1,5 A máximo). Dominio de medida: RN8302B, ISO7741 y acondicionamiento. \\
DVCC\_3v3 & U3: AP63203WU-7 & 3.3 V / 2 A. Dominio digital: ESP32-S3, RS-485, microSD y FGH100M. \\
Respaldo & C16, C17: 2$\times$30 F & Dos supercapacitores de 2,7 V en serie; 15 F equivalentes a 5 V. \\
\bottomrule
\end{tabularx}
\end{table}

\noindent\textbf{Limitador de carga.} La rama de carga del banco de
supercapacitores incorpora un \textbf{Silergy SY6280AAC}, un interruptor de
distribución de 5,5 V y 2 A en SOT-23-5 con límite de corriente programable
\cite{sy6280}. El dispositivo se coloca entre el riel de 5 V de PS1 y el banco
de C16/C17, de modo que la corriente de carga no solicite a PS1 más de lo que
puede entregar. El límite se programa con una resistencia $R_{\mathrm{SET}}$
entre el pin ISET (pin 3) y GND:

\begin{equation}
I_{\mathrm{lim}}\,[\mathrm{A}] = \frac{6800}{R_{\mathrm{SET}}\,[\Omega]}.
\end{equation}

Para un límite inicial de 1 A, $R_{\mathrm{SET}} = 6{,}8\,\mathrm{k\Omega}$;
para 2 A, $R_{\mathrm{SET}} = 3{,}4\,\mathrm{k\Omega}$. El valor definitivo debe
elegirse con el presupuesto de PS1 y la corriente de irrupción admisible. El
SY6280AAC también incorpora protección de sobrecorriente, cortocircuito,
sobretemperatura y bloqueo inverso (sin diodo de cuerpo), lo que evita que el
banco se descargue hacia PS1 cuando el interruptor está apagado.

\noindent\textbf{Ruta de descarga.} La descarga del banco hacia el buck durante
el \emph{Last Gasp} debe definirse explícitamente: el bloqueo inverso del
interruptor impide usar la misma rama de carga como camino de descarga. La red
de D1/D2 debe verificar esa función y, si es necesario, se agrega una ruta de
descarga independiente.

\subsection{Respaldo Last Gasp}

Los supercapacitores C16 y C17 son dos \textbf{HP-2R7-J306VYJ23} de 30 F y
2,7 V. Al conectarlos en serie, la capacidad equivalente es:

\begin{equation}
C_{\mathrm{eq}} = \frac{C_1 C_2}{C_1 + C_2} = \frac{30 \times 30}{30 + 30} = 15\ \mathrm{F}.
\end{equation}

A 5 V, la energía almacenada es:

\begin{equation}
E = \tfrac{1}{2} C_{\mathrm{eq}} V^2 = \tfrac{1}{2}(15)(5)^2 = 187{,}5\ \mathrm{J}.
\end{equation}

El respaldo declarado es de \textbf{3,2 W durante 9 s}, es decir,
$28{,}8\ \mathrm{J}$. Esa cifra representa alrededor del 15\,\pct de la energía
almacenada, por lo que el dimensionamiento es conservador. La autonomía real
depende de la tensión mínima de entrada del buck, de la eficiencia y del
consumo durante el cierre del archivo.

Con el límite de carga del SY6280AAC, el tiempo teórico para llevar el banco de
0 V a 5 V es $t = C_{\mathrm{eq}} \Delta V / I_{\mathrm{lim}}$. Para 1 A,
$t = 15 \times 5 / 1 = 75\ \mathrm{s}$; para 2 A, $37{,}5\ \mathrm{s}$. Son
cotas idealizadas: la corriente decae al final de la carga y el buck consume
parte de la corriente disponible.

El banco incluye resistencias de limitación y balanceo: R27 y R28
($3{,}9\,\mathrm{\Omega}$, 2 W) en serie, R29 y R30 ($1\,\mathrm{k\Omega}$) de
balanceo, y los diodos Schottky D1 y D2 (MS2A40LWS). La topología de carga y
descarga se revisa junto con el layout, antes de fabricar.

\subsection{Presupuesto de corriente}

\begin{table}[H]
\centering
\caption{Presupuesto de corriente del riel de 3.3 V.}
\label{tab:presupuesto}
\small
\begin{tabular}{@{}l r@{}}
\toprule
\textbf{Carga} & \textbf{Consumo (mA)} \\
\midrule
ESP32-S3-WROOM & 500 \\
Quectel FGH100M (Wi-Fi HaLow) & 350 \\
microSD & 200 \\
VISOM RSM3485PCT (RS-485) & 90 \\
RN8302B & 10 \\
\midrule
\textbf{Total} & \textbf{1 150} \\
Capacidad del buck AP63203WU-7 & 2 000 \\
\textbf{Margen} & \textbf{850 (42,5 \pct)} \\
\bottomrule
\end{tabular}
\end{table}

\noindent\textbf{Nota.} El RN8302B se alimenta de \texttt{AVCC\_3v3} (PS2), no
de \texttt{DVCC\_3v3}. Por tanto, el consumo del \texttt{DVCC\_3v3} es de
1 140 mA y el del \texttt{AVCC\_3v3} incluye el AFE y el acondicionamiento. El
total combinado de 1 150 mA es una cota conservadora para el buck. La
asignación por riel debe confirmarse con el consumo medido de cada periférico.

\subsection{Árbol de potencia}

\begin{figure}[H]
\centering
\begin{tikzpicture}[
  x=1cm,y=1cm,
  block/.style={draw,rounded corners=2pt,align=center,font=\sffamily\scriptsize,inner sep=3pt,minimum height=0.9cm},
  src/.style={block,fill=red!8,draw=red!55!black,minimum width=2.8cm},
  pwr/.style={block,fill=green!8,draw=green!45!black,minimum width=3.2cm},
  rail/.style={block,fill=green!18,draw=green!45!black,minimum width=3.2cm,font=\sffamily\scriptsize\bfseries},
  load/.style={block,fill=cyan!8,draw=cyan!45!black,minimum width=2.4cm},
  pwra/.style={-{Latex[length=1.8mm]},semithick,draw=green!45!black,
    rounded corners=9pt},
  lab/.style={font=\sffamily\scriptsize,fill=white,inner sep=1pt}
]
% Fuentes
\node[src] (ac1) at (-1.0,10.0) {Phase\_1 / N};
\node[block,fill=orange!10,draw=orange!55!black,minimum width=2.8cm] (fb1) at (-1.0,8.6) {FB1};
\node[pwr] (ps1) at (-1.0,7.2) {PS1 HLK-10M05\\5 V / 2 A};
\node[rail] (v5) at (-1.0,5.8) {5 V};

\node[src] (ac2) at (7.2,10.0) {Phase\_1 / N};
\node[block,fill=orange!10,draw=orange!55!black,minimum width=2.8cm] (fb2) at (7.2,8.6) {FB2};
\node[pwr] (ps2) at (7.2,7.2) {PS2 HLK-5M03B\\3.3 V / 5 W};
\node[rail] (avcc) at (7.2,5.8) {AVCC\_3v3};

% Respaldo y buck
\node[block,fill=orange!10,draw=orange!55!black,minimum width=2.8cm] (lim) at (-4.6,5.8) {SY6280AAC\\limitador de carga};
\node[block,fill=yellow!18,draw=yellow!45!black,minimum width=2.8cm] (sc) at (-4.6,4.2) {2$\times$30 F\\serie + balanceo};
\node[pwr] (buck) at (-1.0,4.2) {U3 AP63203\\3.3 V / 2 A};
\node[rail] (dvcc) at (-1.0,2.8) {DVCC\_3v3};

% Cargas
\node[load] (esp)  at (-4.4,0.0) {ESP32-S3\\500 mA};
\node[load] (hal)  at (-1.6,0.0) {FGH100M\\350 mA};
\node[load] (sd)   at (1.2,0.0)  {microSD\\200 mA};
\node[load] (rs)   at (4.0,0.0)  {RSM3485PCT\\90 mA};
\node[load] (rn)   at (7.2,0.0)  {RN8302B\\10 mA};

% Conexiones
\draw[pwra] (ac1) -- (fb1);
\draw[pwra] (fb1) -- (ps1);
\draw[pwra] (ps1) -- (v5);
\draw[pwra] (v5.west) -- (lim.east);
\draw[pwra] (lim.south) -- (sc.north);
\draw[pwra, dashed] (sc.east) -- (buck.west);
\draw[pwra] (buck) -- (dvcc);

\draw[pwra] (ac2) -- (fb2);
\draw[pwra] (fb2) -- (ps2);
\draw[pwra] (ps2) -- (avcc);

% Bus DVCC: alimenta el dominio digital, no el RN8302B.
% Cada derivación es un codo ortogonal con el mismo radio, centrado en su carga.
\draw[pwra, line width=1.2pt] (dvcc.south) -- ++(0,-1.15) coordinate (bus);
\draw[pwra, line width=1.2pt] (bus) -| (esp.north);
\draw[pwra, line width=1.2pt] (bus) -| (hal.north);
\draw[pwra, line width=1.2pt] (bus) -| (sd.north);
\draw[pwra, line width=1.2pt] (bus) -| (rs.north);
\fill[green!45!black] (bus) circle (1.6pt);
\draw[pwra] (avcc.south) -- (rn.north);

% Last gasp
\begin{scope}[on background layer]
\node[draw=green!45!black, dashed, rounded corners, fit=(sc), inner sep=6pt,
      label={[font=\sffamily\scriptsize\bfseries,color=green!45!black]above:Last Gasp}] {};
\end{scope}

\end{tikzpicture}
\caption{Árbol de potencia. PS1 alimenta el respaldo y el buck digital a través
del limitador de carga; PS2 alimenta el dominio de medida. Las corrientes
corresponden al presupuesto de la Tabla~\ref{tab:presupuesto}. La línea
discontinua entre el banco y el buck indica la ruta de descarga del
\emph{Last Gasp}.}
\label{fig:power}
\end{figure}

\subsection{Integridad de potencia y térmica}

La corriente de irrupción, la disipación del buck y de R27/R28 durante la
carga, el tiempo real de respaldo y la temperatura de los supercapacitores y
del módulo RS-485 se calculan antes de fabricar. El orden de arranque entre
PS1, PS2 y el buck también debe definirse para evitar corrientes inversas
durante el encendido y el corte de red.
\section{Aspectos por desarrollar}
\label{sec:desarrollo}

La arquitectura descrita está definida a nivel de bloques, componentes e
interfaces. Antes de fabricar la PCB y cerrar la lista de materiales quedan los
puntos que se enumeran a continuación. Se reúnen en esta sección para que el
cuerpo de la nota describa el diseño y no su proceso de revisión.

\subsection{Criterios de medida y calibración}

\begin{itemize}
  \item Definir el modelo de transformador de corriente (núcleo partido,
        salida de 333 mV) y su rango de operación.
  \item Completar la relación de los divisores de tensión, la tensión máxima
        por canal y la potencia disipada en las resistencias de atenuación.
  \item Definir la frecuencia de corte de los filtros RC, la simetría de las
        pistas diferenciales, la ubicación de los capacitores de desacople y el
        plano de referencia del AFE.
  \item Establecer el procedimiento de calibración de ganancia y fase contra un
        patrón, y la verificación de deriva para la clase 2.
\end{itemize}

\subsection{Integridad del aislamiento}

\begin{itemize}
  \item Asignar referencias separadas \texttt{AGND} y \texttt{GND\_DIG} a GND1
        y GND2 del ISO7741, y decidir si se unen en un punto y en cuál.
  \item Verificar el \emph{creepage}, el \emph{clearance} y la separación entre
        la barrera y los planos de cobre en el layout.
  \item Definir la longitud máxima del bus SPI, el retorno de corriente, la
        simetría de las pistas y el tratamiento de los pines EN1/EN2.
  \item Decidir si se implementan los dos GPIO de control del AFE; requieren un
        segundo aislador o un dispositivo de más canales.
\end{itemize}

\subsection{Protección y EMC}

\begin{itemize}
  \item Verificar la tensión nominal de la bornera U4 y de los fusibles
        F1--F3 (250 V) frente a 440 V entre fases en 3P3W.
  \item Justificar el fusible de 2 A frente a la recomendación de 1 A/250 VCA
        del fabricante del HLK-5M03.
  \item Definir el neutro sin fusible, la coordinación de protecciones, la
        energía absorbida por los MOV y el poder de corte de los fusibles.
  \item Definir la unión de \texttt{GND} con la tierra de protección (PE) y el
        comportamiento ante una falla de fase a tierra.
  \item Agregar filtros de modo común y protecciones ESD en el conector USB-C;
        verificar la inmunidad y las emisiones contra IEC 61000-4-2/4/5/6 y
        EN 55011 clase B.
\end{itemize}

\subsection{Alimentación y térmica}

\begin{itemize}
  \item Definir $R_{\mathrm{SET}}$ del SY6280AAC y la ruta de descarga del
        banco hacia el buck durante el \emph{Last Gasp}; verificar la red de
        D1/D2.
  \item Calcular la corriente de irrupción, la disipación del buck y de R27/R28
        durante la carga, el tiempo real de respaldo y la temperatura de los
        supercapacitores y del módulo RS-485.
  \item Definir el orden de arranque entre PS1, PS2 y el buck.
  \item Confirmar la asignación de consumo por riel entre \texttt{DVCC\_3v3} y
        \texttt{AVCC\_3v3}.
\end{itemize}

\subsection{Procesamiento y periféricos}

\begin{itemize}
  \item Integrar el módulo FGH100M y su alimentación, y conectar sus interfaces
        de 6 SDIO, 4 SPI, 2 UART y 2 GPIO.
  \item Completar los capacitores de desacople, las resistencias de pull-up y
        las protecciones ESD de U5, U6, U7 y J401.
  \item Generar la lista de materiales definitiva y el layout de la PCB.
\end{itemize}

\subsection{Revisión normativa}

No se ha realizado la revisión de cumplimiento. Los requisitos de referencia
citan IEC 62053-21, IEC 61010-1, IEC 61000-4-30, IEC 61000-4-2/4/5/6 y
EN 55011. La categoría de medida III, el grado de protección IP y la
temperatura de operación deben confirmarse con el gabinete y el montaje final.

\subsection{Resumen de acciones}

\begin{table}[H]
\centering
\caption{Acciones pendientes antes de fabricar la PCB.}
\label{tab:acciones}
\footnotesize
\begin{tabularx}{\textwidth}{@{}>{\bfseries}p{0.7cm} >{\raggedright\arraybackslash}p{4.4cm} >{\raggedright\arraybackslash}X >{\raggedright\arraybackslash}p{1.5cm}@{}}
\toprule
\textbf{\#} & \textbf{Tema} & \textbf{Acción} & \textbf{Prioridad} \\
\midrule
1 & Referencias del ISO7741 & Separar \texttt{AGND} y \texttt{GND\_DIG}; definir su unión y verificar la separación en el layout. & Alta \\
2 & Canales del AFE & Decidir si se agregan los dos GPIO de control y prever un segundo aislador. & Alta \\
3 & Tensión de entrada & Verificar U4 y F1--F3 para 440 V entre fases en 3P3W. & Alta \\
4 & Coordinación de protecciones & Justificar el fusible de 2 A, calcular la energía de los MOV y el poder de corte. & Media \\
5 & Medida de corriente & Seleccionar el TC y definir el punto de operación. & Alta \\
6 & Divisores y filtros & Completar la relación de los divisores, la potencia disipada y la frecuencia de corte. & Alta \\
7 & Calibración & Definir el procedimiento de ganancia y fase y la verificación de deriva. & Alta \\
8 & Respaldo Last Gasp & Definir $R_{\mathrm{SET}}$, la ruta de descarga y verificar D1/D2. & Alta \\
9 & Térmica & Calcular la disipación de U3 y R27/R28, y la temperatura del banco y del RS-485. & Media \\
10 & Wi-Fi HaLow & Integrar el FGH100M y su alimentación. & Alta \\
11 & Periféricos & Completar desacoples, pull-ups y ESD de U5, U6, U7 y J401. & Alta \\
12 & EMC y normativa & Verificar IEC 61000-4-2/4/5/6, EN 55011, IEC 61010-1 y la categoría de medida III. & Alta \\
\bottomrule
\end{tabularx}
\end{table}
\section{Conclusiones}

El Analizador Trifásico IoT tiene una arquitectura coherente para un medidor
trifásico de gestión: entrada protegida, acondicionamiento analógico dedicado,
AFE RN8302B, aislamiento digital, ESP32-S3 y un conjunto de interfaces adecuado
para RS-485, Wi-Fi HaLow y registro local. La alimentación con respaldo de
supercapacitores cubre el requisito de cierre ordenado y el presupuesto de
corriente del riel digital tiene margen.

La arquitectura, sin embargo, todavía no está lista para fabricar. Los puntos
críticos son la separación de las referencias del ISO7741, la tensión nominal
de la bornera y los fusibles, la integración del FGH100M y de los pasivos de
alimentación y protección de los periféricos, y la falta de cálculos de
protección, filtros y calibración. Ninguno de estos puntos es un obstáculo
insalvable; todos deben cerrarse antes de congelar el diseño y el layout. La
Sección~\ref{sec:desarrollo} reúne las acciones correspondientes.

El valor de la nota está en explicar el porqué de cada bloque y en dejar
explícitos los criterios de selección: protecciones coordinadas, una barrera
galvánica real entre medida y comunicación, y un respaldo dimensionado con
margen sobre la energía necesaria para el cierre ordenado.
\label{sec:referencias}

\addcontentsline{toc}{section}{Referencias}
\begin{thebibliography}{99}

\bibitem{rdem01}
Equipo del proyecto Analizador Trifásico IoT.
\textit{DR-EM-01: Monitor de Energía Industrial Trifásico --- Product and System
Requirements}, Rev. 0.3, octubre de 2026.

\bibitem{diagramav2}
Equipo del proyecto Analizador Trifásico IoT.
\textit{Arquitectura de referencia V2}, 07/10/2026.

\bibitem{rn8302b_ds}
RENERGY.
\textit{RN8302B Datasheet}, versión 3.

\bibitem{rn8302b_um}
RENERGY.
\textit{RN8302B User Manual}, v1.1.

\bibitem{iso7741}
Texas Instruments.
\textit{ISO7740, ISO7741, ISO7742: High-Speed, Robust-EMC Reinforced and Basic
Quad-Channel Digital Isolators}, SLLSEP4G, rev. febrero de 2020.

\bibitem{esp32s3_ds}
Espressif Systems.
\textit{ESP32-S3 Series Datasheet}.

\bibitem{esp32s3_trm}
Espressif Systems.
\textit{ESP32-S3 Technical Reference Manual}.

\bibitem{esp32s3_module}
Espressif Systems.
\textit{ESP32-S3-WROOM-1/1U Datasheet}.

\bibitem{supercap}
Fabricante del supercapacitor.
\textit{Supercapacitor HP-2R7-J306VYJ23 Datasheet}.

\bibitem{ti_duf27a}
Texas Instruments.
\textit{Design Guide: TIDA-010244 --- Three-Phase Shunt-Based Energy Metrology
Reference Design}, literatura TIDUF27A.

\bibitem{ti_duf27a_layout}
Texas Instruments.
\textit{TIDA-010244 Layout Guidelines}.

\bibitem{eval_ade9430}
Analog Devices.
\textit{EVAL-ADE9430: Schematic, Layout and User Guide}.

\bibitem{fgh100m}
Quectel / Morse Micro.
\textit{Wi-Fi HaLow FGH100M EXT01 V30: Reference Design}.

\bibitem{hx_conn}
HX.
\textit{PJ-320A-3P DIP: conector de 3 pines para entrada de corriente}.

\bibitem{jxtconn}
JXTCONN.
\textit{C50126V-4P0G56: bornera de 4 vías, 5 mm, 18 A, 250 V}.
Ficha de producto LCSC, consultada en octubre de 2026.
\url{https://www.lcsc.com/product-detail/C49291885.html}.

\bibitem{jfcfuse}
JDT Fuse.
\textit{JFC2410-0200TS: fusible 2410, 2 A, time-lag}.
Ficha de producto LCSC, consultada en octubre de 2026.
\url{https://www.lcsc.com/product-detail/C55058769.html}.

\bibitem{sy6280}
Silergy.
\textit{SY6280/SY6280A: 5.5 V, 2 A Low Loss Power Distribution Switch},
Rev. 0.3.
\url{https://www.olimex.com/Products/Components/IC/SY6280/resources/SY6280AAC.PDF}.

\bibitem{iec62053}
International Electrotechnical Commission.
\textit{IEC 62053-21: Electricity metering equipment --- Particular
requirements --- Part 21: Static meters for AC active energy (classes 1 and 2)},
2020.

\bibitem{iec61010}
International Electrotechnical Commission.
\textit{IEC 61010-1: Safety requirements for electrical equipment for
measurement, control, and laboratory use --- Part 1: General requirements},
2010.

\bibitem{iec61000430}
International Electrotechnical Commission.
\textit{IEC 61000-4-30: Electromagnetic compatibility (EMC) --- Part 4-30:
Testing and measurement techniques --- Power quality measurement methods},
2015.

\end{thebibliography}

\vspace{0.5em}
\noindent{\footnotesize\color{unalgray}
Las referencias normativas se citan como contexto de los requisitos de
producto; no implican que el prototipo haya sido ensayado ni certificado. Los
enlaces web se consultaron en octubre de 2026.}

\end{document}
